\begin{center}
    \textbf{РЕФЕРАТ}
    \vspace{-12pt}
\end{center}

Расчётно-пояснительная записка содержит содержит: \pageref{bib:end} с., {\the\numexpr\getpagerefnumber{cnt:lists-start}+10} с. с приложениями, \total{figure} рис., \total{table} табл., \total{lstlisting} листинг., \total{applications} прил. и \total{citecnt} источн.

\noindent ГРАФОВЫЕ БАЗЫ ДАННЫХ, РАСПРЕДЕЛЁННЫЕ СИСТЕМЫ, РАСПРЕДЕЛЕНИЕ ВЕРШИН ГРАФА, ШАРДИРОВАНИЕ, ГРАФОВЫЕ АЛГОРИТМЫ, РАСПРЕДЕЛЁННОЕ ХРАНЕНИЕ ДАННЫХ.

Объектом исследования являются распределённые графовые базы данных.

Предметом исследования являются методы и алгоритмы распределения вершин графа между узлами хранения.

Целью работы является разработка программного модуля распределения вершин в графовой базе данных, обеспечивающего снижение количества межшардовых переходов и повышение эффективности обработки графовых запросов.

В ходе работы были выполнены следующие задачи:
\begin{itemize}
\item проанализировать структурные методы оптимизации распределения в графовых базах данных;
\item проанализировать методы потокового распределения в графовых базах данных;
\item проанализировать методы оптимизации распределения в графовых базах данных на основе данных о запросах;
\item определить методы и алгоритмы для дальнейшей реализации;
\item разработать структуры программного обеспечения и определение спецификаций его компонентов;
\item реализоваты компоненты с использованием выбранных средств и их автономное тестирование;
\item оценочно протестировать программное обеспечение, конкретно протестировать кол-во межшардовых переходов при запросах;
\item разработать технологии распределения вершин графа по хранилищам и оптимизация распределения по хранилищам;
\item собрать программного обеспечения и проанализировать его протестировать;
\item разработать технологии запросов на поиск пути к хранилищам и мониторинга их выполнения.
\end{itemize}

В работе выполнен анализ существующих методов распределения графовых данных, включая классические и современные подходы к разбиению графов. Рассмотрены их достоинства и ограничения в контексте применения к распределённым графовым базам данных.

Разработанный программный модуль обеспечивает возможность распределения вершин графа между шардами, анализа качества распределения и проведения экспериментов по оценке эффективности различных алгоритмов. В качестве основной метрики используется количество межшардовых переходов при выполнении запросов поиска пути.

Результаты экспериментов показали, что использование специализированных алгоритмов распределения позволяет существенно снизить количество межшардовых взаимодействий по сравнению со случайным распределением вершин. Это подтверждает эффективность подхода, основанного на учёте структуры графа при его распределении.

Разработанное решение может быть использовано как исследовательский инструмент для анализа алгоритмов распределения графовых данных, а также как основа для дальнейшего развития систем распределённого хранения и обработки графов.
